\subsection{归纳原理}

C61（归纳原理）。设 \(\mathbb{R}\{n\}\) 是理论 \(\mathfrak{C}\) 中的一个关系（其中 \(n\) 不是 \(\mathfrak{C}\) 的常数）。假设关系

\[
\mathrm{R} \{0 \} \text { 且 } (\forall n) ((n \text { 是整数   且 } \mathrm{R} \{n \}) \Longrightarrow \mathrm{R} \{n + 1 \})
\]

在 \(\mathfrak{C}\) 中是定理。在这些条件下，关系

\[
(\forall n) ((n \text {   是整数   }) \Longrightarrow R \{n \})
\]

在 \(\mathfrak{C}\) 中是定理.

我们将用反证法论证。假设关系

\[
(\exists n) (n \text {   是整数   且   } (\text { 非   } R \{n \}))
\]

成立。设 \(q\) 是一个整数，使得“非 \(\mathbb{R}\{q\}\) ”（辅助常数法；参见第一章§3第3号及§4第1号）。整数 \(n\) （满足“ \(n \leqslant q\) 以及（非 \(\mathbb{R}\{n\}\) )''的）构成一个良序非空集合（§3，第2号，注记），因此它有一个最小元素 \(s\) 。如果 \(s = 0\) ，则“非 \(\mathbb{R}\{0\}\) ''，与假设矛盾。若 \(s > 0\) ，那么 \(s = s' + 1\) ，其中 \(s'\) 是一个满足 \(s' < s\) 的整数（第2号，命题2）。根据 \(s\) 的定义，我们有 \(\mathbb{R}\{s'\}\) ，但随后假设蕴含 \(\mathbb{R}\{s\}\) 是真的，这与 \(s\) 的定义相矛盾.

为了应用归纳原理，特别需要证明关系

\[
(n \text {   是整数   且   } R \{n \}) \Longrightarrow R \{n + 1 \}.
\]

为此目的，通常使用辅助假设的方法（第一章，§3，第3号），正是因为这个原因，关系“n是整数且R\{n\}''（甚至R\{n\}）被称为归纳假设。

注记。有许多不同的准则经常以“归纳原理”之名被使用。它们都可以很容易地从C61推导出来，我们在此指出其中最重要的几个：

\begin{enumerate}
\def\labelenumi{(\arabic{enumi})}
\tightlist
\item
  让 \(\mathbf{S}\{n\}\) 为关系
\end{enumerate}

\((\forall p)((n \text{ 是整数 且 } p \text{ 是整数 且 } p < n) \Longrightarrow \mathbb{R}\{p\})\)

并且假设 \(S\{n\}\) 蕴含 \(R\{n\}\) 。则关系

\[
(\forall n) ((n \text {   是整数   }) \Longrightarrow R \{n \})
\]

是真的。因为关系 \(S\{0\}\) 是真的，并且根据假设 \(S\{n\}\) 蕴含 \(R\{n\}\) ；由于关系 \(m < n + 1\) 等价于 \(m \leqslant n\) （第2号，命题2），关系 \(S\{n + 1\}\) 等价于“ \(S\{n\}\) 并且 \(R\{n\}\) ''，因此， \(S\{n\}\) 蕴含 \(S\{n + 1\}\) 。因此，准则C61证明了关系

\[
(\forall n) ((n \text {   是整数   }) \Longrightarrow \mathrm{S} \{n \})
\]

是真的，并且由于 \(\mathbf{S}\{n\}\) 蕴含 \(\mathbf{R}\{n\}\) ，关系

\[
(\forall n) ((n \text {   是整数   }) \Longrightarrow R \{n \})
\]

为真。

\begin{enumerate}
\def\labelenumi{(\arabic{enumi})}
\setcounter{enumi}{1}
\tightlist
\item
  设 \(k\) 为一个整数，设 \(\mathbf{R}\{n\}\) 是一个关系，使得关系
\end{enumerate}

\[
\mathrm{R} \{k \} \text { 且 } (\forall n) ((n \text { 是整数 } \geqslant k \text { 且 } \mathrm{R} \{n \}) \Longrightarrow \mathrm{R} \{n + 1 \})
\]

成立。那么关系

\[
(\forall n) ((n \text {   是整数   } \geqslant k) \Longrightarrow \mathrm{R} \{n \})
\]

成立（“从 \(k\) 开始的归纳”）。令 \(S\{n\}\) 为关系

\[
(n \geqslant k) \Longrightarrow \mathrm{R} \{n \}.
\]

然后通过分情况讨论的方法，我们看到S\{0\} 成立。另一方面，容易验证关系

\[
(n \text {   是整数   且   } S \{n \}) \Longrightarrow S \{n + 1 \}
\]

成立。由C61可知关系

\[
(n \text {   是整数) } \implies \mathrm{S} \{n \}
\]

成立，这证明了我们的断言。

\begin{enumerate}
\def\labelenumi{(\arabic{enumi})}
\setcounter{enumi}{2}
\tightlist
\item
  让 \(a\) 与 \(b\) 是两个整数，使得 \(a \leqslant b\) ，并且让 \(R\{n\}\) 是一个满足
\end{enumerate}

\[
\mathrm{R} \left\{a \right\} \text {   且   } (\forall n) ((n \text {   是整数   且   } a \leqslant n <   b \text {   且   } \mathrm{R} \left\{n \right\}) \Longrightarrow \mathrm{R} \left\{n + 1 \right\}).
\]

那么关系

\[
(\forall n) ((n \text {   是整数   且   } a \leqslant n \leqslant b) \Longrightarrow \mathrm{R} \{n \})
\]

成立。证明与前一情形类似；我们取 S\{n\} 为关系“(a ≤ n \textless{} b) ⇒ R\{n\}''（“限制在区间上的归纳”）。

\begin{enumerate}
\def\labelenumi{(\arabic{enumi})}
\setcounter{enumi}{3}
\tightlist
\item
  让 \(a, b\) 是两个整数，使得 \(a \leqslant b\) ，并且让 \(\mathbf{R}\{n\}\) 是一个满足
\end{enumerate}

\(R\{b\}\) 并且 \((\forall n)((n \text{ 是整数 且 } a \leqslant n < b \text{ 且 } R\{n + 1\}) \Longrightarrow R\{n\}).\)

那么关系

\[
(\forall n) ((n \text {   是整数   且   } a \leqslant n \leqslant b) \Longrightarrow \mathrm{R} \{n \})
\]

是真的。因为我们有关系

（n 是一个整数且 \(a \leqslant n < b\) 以及（非 \(\mathbf{R}\{n\}\) ) \(\Longrightarrow\) 不 \(\mathbf{R}\{n + 1\}\) .

如果对于某个 \(n\) 使得 \(a \leqslant n \leqslant b\) 我们有（并非 \(\mathbf{R}\{n\}\) ），则由(3)可推出（并非 \(\mathbf{R}\{b\}\) ），这与假设矛盾，因此结论成立（“反向归纳”）。

